% Discussion section — DESI paper
% Target: Nature Genetics
% Covers: interpretation, relation to prior work (FitCoal/Cobraa), limitations, future work
% ~700+ words, ≥8 citations

\section*{Discussion}

Our results establish that within-group mean pairwise TMRCA differs markedly and systematically across human continental populations — African populations show $\sim$1.4-fold deeper mean genealogical depth than East Asians across the entire genome. This pattern is incompatible with any panmictic ancestral model, including the severe bottleneck inferred by FitCoal \citep{Hu2023}. It is instead consistent with the ancestral population structure independently inferred by Cobraa \citep{Cousins2025}, providing the first genealogical-depth-based line of evidence corroborating that conclusion. The result extends a long-established body of work showing that African populations harbour greater genomic diversity than non-African populations \citep{Prugnolle2005,DeGiorgio2009}, by demonstrating that this diversity difference extends deep into the Middle Pleistocene as a consequence of ancestral structure rather than post-OoA drift alone.

\subsection*{Relationship to FitCoal and the 930 ka bottleneck}

The FitCoal inference \citep{Hu2023} is based on the site frequency spectrum under the assumption of a panmictic ancestral population. Our simulations confirm that under this framework, all modern continental groups would share the same ancestral genealogy and therefore exhibit equal within-group TMRCA regardless of bottleneck severity or timing. The 370.7 ka AFR$-$EAS difference we observe is thus not a refinement or an extension of the FitCoal finding — it is \textit{inconsistent} with the model class on which FitCoal operates. An SFS-based method cannot distinguish a panmictic bottleneck from ancestral population structure, because both scenarios depress intermediate-frequency alleles and elevate rare alleles \citep{Cousins2025b} — a theoretical degeneracy that has limited the resolution of demographic inference from the SFS alone \citep{Myers2008}. Our TMRCA comparison sidesteps this confounding: under panmixia, AFR and EAS TMRCA are equal regardless of bottleneck magnitude; under structure, they diverge in proportion to the depth and extent of ancestral isolation.

We emphasise that this conclusion does not imply the human ancestral population was large or unconstrained. The two scenarios — panmictic severe bottleneck and structured ancestry with modest deme sizes — are compatible with similar overall genomic diversity. What the DESI results argue against is the specific interpretation that the $\sim$930 ka SFS signal reflects a panmictic reduction in population size. The model-free nature of our test makes it robust to the specific mechanism producing the depth difference: any model that posits a single panmictic ancestral genealogy predicts zero AFR$-$EAS TMRCA difference, and is therefore rejected by our data.

\subsection*{Relationship to Cobraa and ancestral population structure}

The Cobraa structured coalescent model \citep{Cousins2025} infers two ancestral populations diverging at $\sim$1.5 Ma with $\sim$20\% gene flow at $\sim$300 ka. Our Model D simulation, parameterized to approximate this structure, predicts an AFR$-$EAS mean TMRCA difference of 370.1 ka, matching the observed 370.7 ka within measurement uncertainty ($<0.2\%$). This agreement is noteworthy because DESI uses an entirely different data summary (per-window aggregate TMRCA) and a simpler inferential framework than Cobraa (parametric structured HMM). The convergence of two methodologically independent approaches on the same conclusion — that modern human genomic diversity reflects partitioned ancestral lineages rather than a panmictic bottleneck — substantially strengthens both findings.

The biological implication is that approximately 300 ka before the Out-of-Africa dispersal, the proto-Africa and proto-OoA lineages mixed. African modern humans therefore carry genomic segments tracing to both ancestral lineages, contributing the deeper and more heterogeneous within-group TMRCA distribution we observe (AFR bimodality: $\Delta$BIC = $-11{,}170$). Non-African populations exited Africa primarily through the proto-OoA lineage, retaining only the shallower ancestral component and leaving a minimum within-EAS TMRCA component at $\sim$630 ka — far deeper than the $\sim$55 ka expected from the OoA bottleneck itself, and fully consistent with the proto-OoA population having its own deep pre-admixture history.

\subsection*{Limitations and caveats}

Several limitations should be noted. First, DESI does not provide a full demographic model: it tests the panmixia null and quantifies the magnitude of the TMRCA difference, but does not infer split times, admixture proportions, or deme sizes. For explicit demographic parameters, methods such as Cobraa remain more appropriate. Second, absolute TMRCA calibration depends on the assumed mutation rate ($\muup = 1.2 \times 10^{-8}$ per bp per generation); however, our key conclusion — that AFR$-$EAS TMRCA difference is orders of magnitude larger than any panmictic prediction — is not sensitive to $\muup$ choice, as both the observed difference and the panmictic prediction scale identically with $1/\muup$. Third, we did not perform a direct comparison of DESI, FitCoal, and Cobraa applied to the same samples (due to computational constraints), relying instead on simulation-based model comparison. Such a comparison, particularly running FitCoal \citep{Hu2023} and PSMC \citep{Li2011} on the same 1kGP data and contrasting their $\Ne(t)$ trajectories with the DESI depth differences, would provide a more complete picture and remains an important direction for future work. While our simulation results are consistent across three replicates and three models, direct benchmarking on empirical data would further strengthen the conclusions. Fourth, archaic introgression (approximately 2\% Neanderthal DNA in non-African populations) could in principle affect within-EAS TMRCA estimates. Archaic haplotype segments have deeper coalescence times (hundreds of ka to $>$1 Ma), which would if anything increase within-EAS TMRCA — making our reported AFR$>$EAS difference slightly conservative. Fifth, our use of super-population labels (AFR, EAS) rather than individual-level ancestry may mask fine-scale within-group structure; the robustness of the result across 22 autosomes (Z=73) and its replication in AMR sub-populations with known ancestry argue against this being a material concern. Sixth, the per-window K-mixture optimal components (K$>$1) could in part reflect recombination rate heterogeneity across the genome: low-recombination regions retain longer genealogical blocks with more extreme TMRCA values. However, recombination heterogeneity would affect all super-populations similarly; the greater AFR bimodality ($\Delta$BIC = $-11{,}170$) compared to EAS ($\Delta$BIC = $-5{,}875$) is asymmetric in the direction expected from structure, not from recombination variation alone. Seventh, recent rapid population expansion in certain African groups could produce an excess of recent coalescence that might be expected to \textit{reduce} within-AFR TMRCA — not increase it. The deeper AFR genealogies we observe are thus not attributable to recent demography, consistent with large-scale genomic comparisons establishing deep and pervasive African diversity that traces to the Pleistocene \citep{Malaspinas2016,Mallick2016}.

\subsection*{Implications and future directions}

The DESI framework opens several avenues for further investigation. Applying the same within-group TMRCA comparison to ancient DNA samples from Neanderthals, Denisovans, and early modern humans would provide a direct test of whether the proto-Africa and proto-OoA lineages are traceable in the archaic hominin record. Extending the approach to within-Africa continental comparisons (e.g., East African vs. West African vs. southern African populations) could resolve the geographic origin and spatial extent of the deeper ancestral component \citep{Cousins2025}. Integration with ancient DNA from the African Middle Stone Age, now being recovered in increasing numbers \citep{Stringer2016}, would provide direct calibration of the ancestral structure signal. More broadly, genealogical depth comparisons across populations represent a largely unexploited dimension of demographic inference: any scenario that produces asymmetric genealogical depth between groups — archaic introgression, local adaptation, sex-biased dispersal — leaves a detectable signature in within-group TMRCA. \desi{} provides a computationally tractable and statistically principled framework for exploiting this signal at genome scale.

The results presented here illustrate a general principle: that the interpretation of an observed genomic signal depends critically on the demographic model assumed. When the panmixia assumption is relaxed and within-group genealogical depth is compared directly, the $\sim$930 ka signal in human genomic data is better described as a signature of ancestral population structure than as evidence for a catastrophic panmictic bottleneck. Whether this structure reflects two distinct proto-human populations, an extended contact zone, or a more complex network of interconnected demes remains an open question — one that will require integration of genomic, fossil, and archaeological evidence to resolve.
